\documentclass[../DoAn.tex]{subfiles}
\begin{document}

\begin{center}
    \Large{\textbf{TÓM TẮT NỘI DUNG BÁO CÁO NGHIÊN CỨU TỐT NGHIỆP 2}}\\
\end{center}
\vspace{1cm}

Trong kỷ nguyên số hóa, nhu cầu học tiếng Nhật ngày càng tăng cao nhưng người học thường gặp khó khăn do lộ trình học rời rạc và thiếu công cụ tra cứu ngữ cảnh trực quan. Các hệ thống học tập trực tuyến (LMS) hiện nay đa số tập trung vào truyền tải bài giảng video tĩnh, thiếu tính năng ôn tập từ vựng chủ động như thẻ flashcard hoặc các trợ lý AI thông minh để giải thích ngữ pháp cá nhân hóa theo thời gian thực. Để giải quyết các hạn chế này, em đã chọn hướng tiếp cận xây dựng một nền tảng e-learning tích hợp toàn diện (All-in-One), kết hợp quản lý học liệu LMS truyền thống với các công cụ tương tác thông minh như ôn tập lật thẻ flashcard, tra cứu từ điển popup nhúng trực tiếp và trợ lý học tập tích hợp Google Gemini AI. Việc lựa chọn giải pháp này xuất phát từ mong muốn tối ưu hóa trải nghiệm học tập, giúp học viên không bị ngắt quãng khi phải sử dụng nhiều ứng dụng khác nhau để học lý thuyết, tra từ điển và luyện thi thử.

Hệ thống JPMaster được phát triển dựa trên kiến trúc phân tách rõ ràng giữa Frontend SPA (React, Vite, Framer Motion) và Backend API (Express, TypeScript, PostgreSQL), đồng thời tích hợp các dịch vụ bên ngoài như PayOS cho việc thanh toán tự động và Cloudinary để quản lý tài nguyên đa phương tiện. Đóng góp chính của báo cáo là thiết kế và hiện thực hóa thành công một hệ thống học tập khép kín: từ việc mua khóa học qua cổng thanh toán tự động, theo dõi tiến độ học chi tiết, ôn tập từ vựng thông minh đến làm các đề thi thử JLPT theo cấu trúc chuẩn và nhận sự trợ giúp trực tiếp của AI theo ngữ cảnh bài học. Kết quả thu được là một sản phẩm phần mềm hoàn chỉnh chạy thử nghiệm ổn định, vượt qua các ca kiểm thử tự động và thủ công, sẵn sàng hỗ trợ hiệu quả cho người học tiếng Nhật trên mọi thiết bị.

\vspace{1cm}
\begin{flushright}
Sinh viên thực hiện\\
\begin{tabular}{@{}c@{}}
\textit{(Ký và ghi rõ họ tên)}
\end{tabular}
\end{flushright}

\end{document}